\documentclass{beamer}
% \usepackage[utf8]{inputenc}
\usepackage[T1]{fontenc}


\usepackage[ngerman]{babel}

\usepackage{tikz}
\usetikzlibrary{decorations,shapes,positioning,fit}



\title{Cryptoparty: Einführung}
\subtitle{Eine Einführung in E-Mail-Sicherheit mit GPG}
\author{iFSR TU Dresden}

\newcommand{\alice}[1][1cm]{\includegraphics[width=#1]{img-src/alice.pdf}}
\newcommand{\bob}[1][1cm]{\includegraphics[width=#1]{img-src/bob.pdf}}

\begin{document}

\begin{frame}
  \maketitle
\end{frame}

\begin{frame}
  \frametitle{Zum Verlauf der Veranstaltung}
  \framesubtitle{oder: Willkommen!}
  \begin{block}{Dreiteilige Veranstaltung}
  \begin{enumerate}
  \item Zuerst: Konzeptuelle Einführung in E-Mail-Sicherheit mit GPG (Schlüssel
    und Vertrauen)
  \item Betriebssystemspezifische Einführungen in die Verwendung von GPG um
    E-Mails zu verschlüsseln
  \item Gegenseitiger Austausch von Schlüsseln austauschen und gegenseitiges
    Signieren
  \end{enumerate}
  \end{block}
\end{frame}

\newcommand{\pollq}[1]{\structure{#1}}

\begin{frame}
  \frametitle{Publikumsumfragen}
  \only<1>{\pollq{Wer hat schon einmal\vphantom{g} via E-Mail mit Behörden kommuniziert?}}
  \only<2>{\pollq{Wer hat schon einmal\vphantom{g} Kontodaten via E-Mail versendet?}}
  \only<3>{\pollq{Wer hat schon einmal eine Rechnung nur via E-Mail erhalten?}}
\end{frame}

\begin{frame}
  \frametitle{Das Konzept für E-Mail-Sicherheit mit GPG}
  \begin{block}{Übertragung}
    \begin{center}
      \begin{tikzpicture}
        \path (0,0) node (alice) {\alice};
        \path (1.5,0) node (gpg-alice) {GPG};
        \path (8.5,0) node (gpg-bob) {GPG};
        \path (10,0) node (bob) {\bob};


        \draw[-latex] (alice) -- (gpg-alice);
        \draw[-latex,double] (gpg-alice) -- (gpg-bob);
        \draw[-latex] (gpg-bob) -- (bob);

        \path (5,0) node [draw,fill=white, cloud, cloud puffs=14, minimum width=3.5cm, minimum height=2cm] (internet) {};
        \path (internet) node[yshift=0.25cm] {Internet};
      \end{tikzpicture}
    \end{center}
    \begin{itemize}
    \item GPG bei Alice: verschlüsselt und signiert
    \item GPG bei Bob: entschlüsselt und prüft Signatur
    \end{itemize}
  \end{block}
\end{frame}

\begin{frame}
  \frametitle{Schlüssel bei GPG}
  \begin{columns}
    \column{0.3\textwidth}
    \begin{center}
      \usebeamercolor{frametitle}
      \begin{tikzpicture}[every node/.style={inner sep=0.2cm}]
        \node (pub) {\texttt{pub}};
        \node [right=of pub] (sec) {\texttt{sec}};
        \node [above=of pub,yshift=-0.8cm] (label) {Schlüssel};

        \node [fill=black!15!white,fit=(pub)(sec)(label),draw] (key) {};

        % re-painting the nodes as they got visually killed by the filled fit node
        \path (pub) node [fill=white,draw,rectangle,dashed] {\texttt{pub}};
        \path (sec) node {\texttt{sec}};
        \path (label) node [color=fg] {Schlüssel};
      \end{tikzpicture}
    \end{center}
    \column{0.65\textwidth}
    \begin{itemize}
    \item Schlüsselpaar
    \item Öffentlicher Teil (\texttt{pub}): \\zum Verschlüsseln \\und Prüfen von Signaturen
    \item Geheimer Teil (\texttt{sec}): \\zum Entschlüsseln \\und Erstellen von Signaturen
    \item Zusätzlich: Informationen \\z.B. Name, E-Mail-Adresse
    \end{itemize}
  \end{columns}
\end{frame}

\begin{frame}
  \frametitle{``Ich habe doch nichts zu verbergen''}
  \framesubtitle{oder: warum man E-Mails verschlüsseln möchte und warum es trotzdem nichts bringt}

  \begin{block}{Warum?}
    \begin{itemize}
    \item Jeder hat etwas zu verbergen (Kontonummer? Konto\emph{stand}?
      Rechnungen? Beziehungen?)
    \item Daten und Text werden oft unverschlüsselt\\via E-Mail übertragen
    \small
    \item Praktisch jeder kann lesen was unverschlüsselt per E-Mail übertragen
      wird
    \item E-Mails sind leicht zu fälschen
    \end{itemize}
  \end{block}
  \pause
  \begin{block}{Probleme}
    \begin{itemize}
    \item Selbst mit GPG wird nicht alles verschlüsselt
    \item Gegenüber muss GPG verwenden
    \end{itemize}
  \end{block}
\end{frame}

\begin{frame}
  \frametitle{Was ist (nicht) abgedeckt?}
  \begin{columns}
    \column[t]{0.45\textwidth}
    \begin{block}{Sicher}
      \begin{itemize}
      \item Text
      \item ggfs. Anhänge
      \end{itemize}
    \end{block}
    \column[t]{0.45\textwidth}
    \begin{block}{Unsicher}
      \begin{itemize}
      \item Absender (teilw.)
      \item Empfänger
      \item Datum
      \item Betreff
      \item und mehr
      \end{itemize}
    \end{block}
  \end{columns}
\end{frame}

\begin{frame}
  \frametitle{Über Verschlüsselung hinaus: Signaturen}
  GPG bietet auch Signaturen, digitale Unterschriften
  \begin{itemize}
  \item Deckt Text und ggfs. Anhänge ab
  \item Manipulationen an Text und Anhängen sind erkennbar
  \item Unterschrift ist mit dem Schlüssel assoziiert
  \end{itemize}
\end{frame}

\end{document}
